\documentclass[11pt]{article}

% Packages
\usepackage[utf8]{inputenc}
\usepackage[margin=1in]{geometry}
\usepackage{amsmath, amssymb}
\usepackage{graphicx}

% Document info
\title{Deep Learning - Assignment 1}
\author{Sebastian Nickels - snickels3@gatech.edu}

\begin{document}

\maketitle

\section{Theory Problem Set}

\newpage

\subsection*{1. Problem}
We have the softmax function $$s_i(z) = \frac{e^{z_i}}{\sum_k e^{z_k}}$$

We want to know $$\frac{\partial s}{\partial z}$$

For that, we will look at $$\frac{\partial s_i}{\partial z}$$

Which consists of $$\frac{\partial s_i}{\partial z_j}$$

$$\frac{\partial s_i}{\partial z_i} = \frac{\partial }{\partial z_i} \left(\frac{e^{z_i}}{\sum_k e^{z_k}}\right)$$

Set $G = \sum_k e^{z_k}$

There are two cases:

First case: $i = j$:

\begin{align*}
\frac{\partial s_i}{\partial z_i} &= \frac{\partial }{\partial z_i} (\frac{e^{z_i}}{G}) \\
&= \frac{e^{z_i} G - e^{z_i} e^{z_i}}{G^2} & \text{(quotient rule)}  \\
&= \frac{e^{z_i}(G - e^{z_i})}{G^2} \\
&= \frac{e^{z_i}}{G^2}(G - e^{z_i}) \\
&= \frac{e^{z_i}}{G^2} \frac{G}{G} (G - e^{z_i}) \\
&= \frac{e^{z_i}}{G^2} G (1 - \frac{e^{z_i}}{G}) \\
&= \frac{e^{z_i}}{G} (1 - \frac{e^{z_i}}{G}) \\
&= s_i (1 - s_i) \\
\end{align*}

Second case: $i \neq j$:
\begin{align*}
\frac{\partial s_i}{\partial z_i} &= \frac{\partial }{\partial z_i} (\frac{e^{z_i}}{G}) \\
&= \frac{0 e^{z_j} - e^{z_i} e^{z_j}}{G^2} & \text{(quotient rule)} \\
&= - \frac{e^{z_i} e^{z_j}}{G^2} \\
&= - \frac{e^{z_i}}{G} \frac{e^{z_j}}{G} \\
&= - s_i s_j \\
\end{align*}

These two cases can be combined using the Kronecker delta $\delta_{ij}$:
$$\frac{\partial s_i}{\partial z_j} = s_i (\delta_{ij} - s_j)$$

Which leads to the result:
$$
\frac{\partial \mathbf{s}}{\partial \mathbf{z}}
=
\begin{bmatrix}
s_1(1 - s_1) & -s_1 s_2 & \cdots & -s_1 s_n \\
-s_2 s_1 & s_2(1 - s_2) & \cdots & -s_2 s_n \\
\vdots & \vdots & \ddots & \vdots \\
-s_n s_1 & -s_n s_2 & \cdots & s_n(1 - s_n)
\end{bmatrix}
$$

\newpage

\subsection*{2. Problem}
$$
w_{\text{AND}} =
\begin{pmatrix}
1 \\
1
\end{pmatrix}
$$

$$
b_{\text{AND}} = -1.5
$$

$$
w_{\text{OR}} =
\begin{pmatrix}
1 \\
1
\end{pmatrix}
$$
$$
b_{\text{OR}} = -0.5
$$

I chose the same weight vector for both, but basically just change the bias depending on the function, as f = 1 for AND needs basically more positive values and is a subset of OR.

\newpage

\subsection*{3. Problem}
Let's prove this by trying to solve the linear equation system:
1. 
\begin{align*}
x = (0, 0) &\Rightarrow w_1 \cdot 0 + w_2 \cdot 0 + b < 0 &\text{as } f(x) = 0 \\
&\Rightarrow b < 0 \\
\end{align*}

2.
\begin{align*}
x = (0, 1) &\Rightarrow w_1 \cdot 0 + w_2 \cdot 1 + b \geq 0 &\text{as } f(x) = 1 \\
&\Rightarrow w_2 + b \geq 0 \\
&\Rightarrow w_2 \geq -b \\
\end{align*}

3.
\begin{align*}
x = (1, 0) &\Rightarrow w_1 \cdot 1 + w_2 \cdot 0 + b \geq 0 &\text{as } f(x) = 1 \\
&\Rightarrow w_1 + b \geq 0 \\
&\Rightarrow w_1 \geq -b \\
\end{align*}

4.
\begin{align*}
x = (1, 1) &\Rightarrow w_1 \cdot 1 + w_2 \cdot 1 + b < 0 &\text{as } f(x) = 0 \\
&\Rightarrow w_1 + w_2 + b < 0 \\
&\Rightarrow w_1 + w_2 < -b \\
\end{align*}

Take 2 and 3
\begin{align*}
w_2 \geq -b \land w_1 \geq -b &\Rightarrow w_1 + w_2 \geq -b + -b \\
&\Rightarrow w_1 + w_2 \geq -2b \
\end{align*}

This is incompatible with 4, as it would require

$$w_1 + w_2 \ge -2b > -b > w_1 + w_2$$

As we have $b < 0$, and even without that constraint, we cannot have $w_1 + w_2 > w_1 + w_2$.

Therefore we showed that there XOR can not be represented using a linear model, as there is no linear combination of weights (and bias) that separates them, which can be also seen when plotting them into a graph and trying to draw a line that separates those two classes perfectly (it won't work).

\newpage

\section{Paper Review}
For this section, I reviewed the paper "Understanding deep learning requires rethinking generalization" by Chiyuan Zhang, Samy Bengio, Moritz Hardt, Benjamin Recht, Oriol Vinyals.

Short review:

In this section, you must choose one of the papers below and complete the
following:
1. provide a short review of the paper,
2. answer paper-specific questions,
Guidelines: Please restrict your reviews to no more than 350 words and
answers to questions to no more than 350 words per question. The review
part (1) should include answers to the following:
3. What is the main contribution of this paper? In other words, briefly
summarize its key insights. What are some strengths and weaknesses
of this paper?
4. What is your personal takeaway from this paper? This could be ex-
pressed either in terms of your perceived novelty of this paper compared
to others you’ve read in the field, potential future directions of research
in the area that the authors haven’t addressed, or anything else that
struck you as being noteworthy

The second paper is again one that questions conventional wisdom and shows
that large neural networks can actually fit random labels, that is labels that
remain fixed but, for example, have no relationship to what is actually in the
image. The paper title is “Understanding deep learning requires rethinking
generalization” and can be found here.
Questions for this paper:
• If neural networks can “memorize” the data, which is the only thing
they can do for random label assignments that don’t correlate with
patterns in the data, why do you think neural networks learn more
meaningful, generalizable representations when there are meaningful
patterns in the data?
• How does this finding align or not align with your understanding of
machine learning and generalization?

\section{Report}
% WIP

\end{document}
